\SourceLeaf{338}

\begin{center}
[\emph{Vorlesung über die Theorie der imaginären Grössen.}]

Vom Viereck.

[1.]
\end{center}

Hat man ein Viereck (1), (2), (3), (4), das durch die complexen Zahlen $p'$, $p''$, $p'''$,
$p^{\mathrm{IV}}$ bestimmt ist, so lässt sich eine allgemeine Formel vermittelst der Lateralen
bestimmen, indem man auf ein Dreieck (I), (II), (III) recurrirt, das daraus auf die alleinig
symmetrische Art entstanden ist, wo also die Eckpunkte folgendermassen festgelegt sind:
\[
\begin{array}{rll}
\text{(I)} & \text{durch} & p'p'' + p'''p^{\mathrm{IV}} \\
\text{(II)} & \text{,,} & p'p''' + p'' p^{\mathrm{IV}} \\
\text{(III)} & \text{,,} & p'p^{\mathrm{IV}} + p'' p'''.
\end{array}
\]
Ist z.\,B. die Complexe, die
\[
\begin{array}{rlll}
\text{(1)} & \text{festlegt} & : & 2 + \sqrt{-1} = p', \\
\text{die (2)} & \text{,,} & : & 3 \hspace{2.8em} = p'', \\
\text{,, (3)} & \text{,,} & : & 1 + 2\sqrt{-1} = p''', \\
\text{,, (4)} & \text{,,} & : & 1 + 3\sqrt{-1} = p^{\mathrm{IV}},
\end{array}
\]
so wird
\[
\begin{array}{rlll}
\text{(I)} & \text{bestimmt durch} & : & 1 + 8\sqrt{-1} \\
\text{(II)} & \text{,, ,,} & : & 3 + 14\sqrt{-1} \\
\text{(III)} & \text{,, ,,} & : & 2 + 13\sqrt{-1}.
\end{array}
\]

Beziehe ich (II) auf (I), so erhalte ich, indem ich die Complexen, die zu diesen Ecken gehören,
von einander abziehe:
\[
\text{(II)} - \text{(I)} = p'(p''' - p'') + p^{\mathrm{IV}}(p'' - p'''),
\]
und ebenso:
\[
\text{(III)} - \text{(II)} = p'(p^{\mathrm{IV}} - p''') + p'' (p''' - p^{\mathrm{IV}})
\]
\[
\text{(I)} - \text{(III)} = p'(p'' - p^{\mathrm{IV}}) + p'''(p^{\mathrm{IV}} - p'')
\]
oder
\[
\begin{array}{rl}
\text{(II)} - \text{(I)} & = (p' - p^{\mathrm{IV}})(p''' - p'') \\
\text{(III)} - \text{(II)} & = (p' - p'')(p^{\mathrm{IV}} - p''') \\
\text{(I)} - \text{(III)} & = (p' - p''')(p'' - p^{\mathrm{IV}}).
\end{array}
\]

Beziehe ich die Winkel des Vierecks auf eine feste Axe, und bezeichne z.\,B. den Winkel, den (4,1)
mit dieser bildet, mit [4,1], so ist:
\[
\begin{array}{rl}
-p^{\mathrm{IV}} + p' & = (4,1)\{\cos[4,1] + i\sin[4,1]\} \\
p''' - p'' & = (3,2)\{\cos[3,2] + i\sin[3,2]\}
\end{array}
\]
\[
\text{etc.}
\]
Man hat daher:
\[
\begin{array}{rl}
\text{(I, II)} & = (4,1) . (3,2) \\
\text{(II, III)} & = (2,1) . (4,3) \\
\text{(III, I)} & = (3,1) . (4,2).
\end{array}
\]
